%% The MIT License (MIT)
%%
%% Copyright (c) 2015 Daniil Belyakov
%%
%% Permission is hereby granted, free of charge, to any person obtaining a copy
%% of this software and associated documentation files (the "Software"), to deal
%% in the Software without restriction, including without limitation the rights
%% to use, copy, modify, merge, publish, distribute, sublicense, and/or sell
%% copies of the Software, and to permit persons to whom the Software is
%% furnished to do so, subject to the following conditions:
%%
%% The above copyright notice and this permission notice shall be included in all
%% copies or substantial portions of the Software.
%%
%% THE SOFTWARE IS PROVIDED "AS IS", WITHOUT WARRANTY OF ANY KIND, EXPRESS OR
%% IMPLIED, INCLUDING BUT NOT LIMITED TO THE WARRANTIES OF MERCHANTABILITY,
%% FITNESS FOR A PARTICULAR PURPOSE AND NONINFRINGEMENT. IN NO EVENT SHALL THE
%% AUTHORS OR COPYRIGHT HOLDERS BE LIABLE FOR ANY CLAIM, DAMAGES OR OTHER
%% LIABILITY, WHETHER IN AN ACTION OF CONTRACT, TORT OR OTHERWISE, ARISING FROM,
%% OUT OF OR IN CONNECTION WITH THE SOFTWARE OR THE USE OR OTHER DEALINGS IN THE
%% SOFTWARE.

\documentclass{mcdowellcv}
\usepackage[hidelinks,unicode]{hyperref}
\usepackage{resume-variant}
\hypersetup{
    pdftitle={Tommaso Bonato - Resume},
    pdfauthor={Tommaso Bonato}
}
\urlstyle{same}

\name{Tommaso Bonato}
\contacts{\href{mailto:tommasobonato95@gmail.com}{\nolinkurl{tommasobonato95@gmail.com}}
    \enspace\textbar\enspace
    \href{https://tommasobonato.github.io/}{tommasobonato.github.io}}

\begin{document}
\makeheader

\begin{cvsection}{EDUCATION}
    \begin{cvsubsection}{Ph.D. in Computer Science}{ETH Zürich}{Jan 2023 -- Mar 2027 (expected)}
        \begin{itemize}
            \item Part of the Scalable Parallel Computing Lab (SPCL), advised by Prof. Torsten Hoefler.
            \item Research: AI workloads, collective communication, congestion control, and load balancing for AI/HPC.
            \item Publications in venues including SC, NSDI, SIGCOMM, EuroSys, and \textit{Communications of the ACM}.
            \item Supervised and mentored several bachelor's and master's students, guiding research design and execution.
            \item Teaching Assistant: Wireless Networking and Mobile Computing, Fall 2026.
        \end{itemize}
    \end{cvsubsection}

    \begin{cvsubsection}{M.S. in Computer Science}{ETH Zürich}{Sep 2018 -- Nov 2021}
        \begin{itemize}
            \item GPA: 5.31/6.00. Thesis on adaptive routing in the Slingshot interconnect, with Hewlett-Packard.
        \end{itemize}
    \end{cvsubsection}

    \begin{cvsubsection}{B.S. in Computer Science}{University of Bologna}{Sep 2014 -- Mar 2018}
        \begin{itemize}
            \item Graduated with honors, 110/110; GPA: 29.1/30.
            \item Scholarship-funded exchange at the University of California, Santa Cruz (2016--2017).
        \end{itemize}
    \end{cvsubsection}
\end{cvsection}

\begin{cvsection}{WORK EXPERIENCE}
    \begin{cvsubsection}{HPC and AI Software Network Architect Intern}{NVIDIA}{Jun 2026 -- Nov 2026}
        \begin{itemize}
            \item Design and evaluate multi-NIC load balancing for GPU-initiated networking (GIN) in AI systems.
            \item Contributed to the \highlight{production release} of a multi-NIC load-balancing feature for NVIDIA Vera Rubin.
        \end{itemize}
    \end{cvsubsection}

    \begin{cvsubsection}{Research Engineer (Part-Time)}{Microsoft (via Experis)}{Aug 2022 -- Mar 2026}
        \begin{itemize}
            \item Designed transport protocols for Azure's AI/LLM and HPC infrastructure in the Network Architecture team.
            \item Represented Microsoft in the Ultra Ethernet (UE) Transport Working Group, contributing to congestion control and load balancing.
            \item Developed proofs of concept for inter-datacenter congestion control and reliability, networking/AI co-design, and packet- and flow-level load balancing.
            \item Extended Microsoft's C++ htsim simulator for \highlight{workload-driven evaluation} of AI/HPC networking.
        \end{itemize}
    \end{cvsubsection}

    \begin{cvsubsection}{Software Engineer Intern}{Amazon}{Feb 2019 -- Jul 2019}
        \begin{itemize}
            \item Built an LSTM-based transportation capacity recommendation system using Keras and AWS, contributing to estimated annual savings of \highlight{USD 5 million}.
            \item Built alerts that caught high-priority transportation capacity misconfigurations for Prime Day and holidays.
        \end{itemize}
    \end{cvsubsection}

    \begin{cvsubsection}{Software Engineer Intern}{European Space Agency}{Jul 2017 -- Oct 2017}
        \begin{itemize}
            \item Developed data analysis tools for the James Webb Space Telescope; integrated Sandia National Laboratories' C++ libraries into Fortran code, achieving a \highlight{4\texttimes{} speedup}.
        \end{itemize}
    \end{cvsubsection}
\end{cvsection}

\begin{cvsection}{SELECTED PUBLICATIONS}
    \begin{itemize}[itemsep=4pt]
        \item \href{https://doi.org/10.1145/3767295.3769320}{\textbf{REPS: Recycled Entropy Packet Spraying for Adaptive Load Balancing and Failure Mitigation.}}
            \textit{EuroSys 2026; first author.} Reference load-balancing design for \highlight{Ultra Ethernet (UE)} and inspiration for \highlight{Multipath Reliable Connection (MRC)}.
        \item \href{https://doi.org/10.1145/3712285.3759838}{\textbf{ATLAHS: An Application-centric Network Simulator Toolchain for AI, HPC, and Distributed Storage.}}
            \textit{SC 2025; co-first author.} \award{Best Student Paper Finalist.}
        \item \href{https://doi.org/10.1145/3712285.3759884}{\textbf{Uno: A One-Stop Solution for Inter- and Intra-Data Center Congestion Control and Reliable Connectivity.}}
            \textit{SC 2025; co-first author.}
        \item \href{https://arxiv.org/abs/2508.16809}{\textbf{PICO: Performance Insights for Collective Operations.}}
            \textit{ISC 2026.} \award{Hans Meuer Award (Best Paper).} MPI/NCCL benchmarking and tuning for LLM training.
    \end{itemize}
\end{cvsection}

\begin{cvsection}{TECHNICAL SKILLS}
    \begin{itemize}
        \item \textbf{Proficient:} C/C++, Python. \textbf{Familiar:} Java, MPI, NCCL, RDMA, Bash, Git, Docker, Linux.
        \item \textbf{Expertise:} LLM training/inference workloads, GPU communication. \textbf{Simulation:} htsim, ns-3, SST.
    \end{itemize}
\end{cvsection}
\end{document}
